\documentclass[12pt]{article}\usepackage[]{graphicx}\usepackage[]{color}
%% maxwidth is the original width if it is less than linewidth
%% otherwise use linewidth (to make sure the graphics do not exceed the margin)
\makeatletter
\def\maxwidth{ %
  \ifdim\Gin@nat@width>\linewidth
    \linewidth
  \else
    \Gin@nat@width
  \fi
}
\makeatother

\definecolor{fgcolor}{rgb}{0.345, 0.345, 0.345}
\newcommand{\hlnum}[1]{\textcolor[rgb]{0.686,0.059,0.569}{#1}}%
\newcommand{\hlstr}[1]{\textcolor[rgb]{0.192,0.494,0.8}{#1}}%
\newcommand{\hlcom}[1]{\textcolor[rgb]{0.678,0.584,0.686}{\textit{#1}}}%
\newcommand{\hlopt}[1]{\textcolor[rgb]{0,0,0}{#1}}%
\newcommand{\hlstd}[1]{\textcolor[rgb]{0.345,0.345,0.345}{#1}}%
\newcommand{\hlkwa}[1]{\textcolor[rgb]{0.161,0.373,0.58}{\textbf{#1}}}%
\newcommand{\hlkwb}[1]{\textcolor[rgb]{0.69,0.353,0.396}{#1}}%
\newcommand{\hlkwc}[1]{\textcolor[rgb]{0.333,0.667,0.333}{#1}}%
\newcommand{\hlkwd}[1]{\textcolor[rgb]{0.737,0.353,0.396}{\textbf{#1}}}%
\let\hlipl\hlkwb

\usepackage{framed}
\makeatletter
\newenvironment{kframe}{%
 \def\at@end@of@kframe{}%
 \ifinner\ifhmode%
  \def\at@end@of@kframe{\end{minipage}}%
  \begin{minipage}{\columnwidth}%
 \fi\fi%
 \def\FrameCommand##1{\hskip\@totalleftmargin \hskip-\fboxsep
 \colorbox{shadecolor}{##1}\hskip-\fboxsep
     % There is no \\@totalrightmargin, so:
     \hskip-\linewidth \hskip-\@totalleftmargin \hskip\columnwidth}%
 \MakeFramed {\advance\hsize-\width
   \@totalleftmargin\z@ \linewidth\hsize
   \@setminipage}}%
 {\par\unskip\endMakeFramed%
 \at@end@of@kframe}
\makeatother

\definecolor{shadecolor}{rgb}{.97, .97, .97}
\definecolor{messagecolor}{rgb}{0, 0, 0}
\definecolor{warningcolor}{rgb}{1, 0, 1}
\definecolor{errorcolor}{rgb}{1, 0, 0}
\newenvironment{knitrout}{}{} % an empty environment to be redefined in TeX

\usepackage{alltt}
\usepackage{scrtime} % for \thistime (this package MUST be listed first!)
\usepackage[margin=1in]{geometry}
\usepackage{fancyhdr}
\pagestyle{fancyplain}
\usepackage[T1]{fontenc}
\usepackage{tikz-cd}
\usepackage{amsmath}
% So David can comment:
\newcommand{\de}[1]{{\color{red}{\bfseries DE:} #1}}
% Use the R logo when referring to R:
\usepackage{xspace}
\newcommand{\Rlogo}{\protect\includegraphics[height=2ex,keepaspectratio]
{images/Rlogo.pdf}\xspace} 
% Note: must have the R logo saved in
%       an "images" directory



\rhead{\fancyplain{}{4MB3 Project Notebook 2018 \hfill Daniel Presta}}
\title{Math 4MB3 Project Notebook 2018}
\author{Daniel Presta (The Infective Collective)}
\date{\today\ @ \thistime}
\IfFileExists{upquote.sty}{\usepackage{upquote}}{}
\begin{document}
\maketitle

\section*{Friday 9 March 2018}

\textbf{Group Meeting (in class)} \\
\emph{Approximate Duration: 0.5 Hours}

\begin{itemize}
\item We discussed project topics, and settled on ``Spatial Epidemic Dynamics: Synchronization" as our topic
\item We sent an email to Dr. Earn, confirming our topic choice.
\end{itemize}

\section*{Wednesday 14 March 2018}

\textbf{Group Meeting (in class)} \\
\emph{Approximate Duration: 1  Hour}

\begin{itemize}
\item Updated ``README.md" in the GitHub repository; added papers that should be read for the next meeting.
\item Discuss step (b) (construction of spatial SIR model) next meeting, will hopefully have model approved by Dr. Earn by the end of the next meeting.
\item Read suggested papers.
\end{itemize}

\section*{Friday 16 March 2018}

\textbf{Group Meeting (in class)} \\
\emph{Approximate Duration: 1 Hour}

\begin{itemize}
\item Tried to construct spatial SIR model after reading suggested papers.
\item We approached Dr. Earn about a possible idea for a model.
\end{itemize}

\textbf{Group Meeting} \\
\emph{Approximate Duration: 1 Hour}

\begin{itemize}
\item Split up work 
\item Aurora and Michael are to work on numerical/computational simulations.
\item DPark and DPresta will try to connect this model to existing theorems
\item DPark and DPresta will try to use existing theorems to analytically determine/prove a set of criteria for coherence in our model.
\end{itemize}

\section*{Tuesday 20 March 2018}

\textbf{Solo Work} \\
\emph{Approximate Duration: 2 Hours}

\begin{itemize}
\item Read (Earn \& Levin 2006) and (McCluskey \& Earn 2011) to better understand analytical approaches
\item Using their approaches, attempted to derive a set of criteria for coherence in our model.
\item Messed around with possible reproductive functions, before stopping and realizing that the dispersal matrix in the exponent cannot be simplified.
\end{itemize}

\section*{Wednesday 21 March 2018}

\textbf{Solo Work} \\
\emph{Approximate Duration: 2 Hours}

\begin{itemize}
\item Took different approach to finding analytical condition.
\item Like DPark, I analyzed a two patch SIS model and tried to derive a coherence criterion.
\item Difference is that I did not assume equal coupling.
\item Algebra did not simplify, ran into a wall.
\end{itemize}

\textbf{Group meeting} \\
\emph{Approximate Duration: 1 Hour}
\begin{itemize}
\item Aurora, Michael, and DPark all troubleshooted code and fixed their error in the Rcpp file.
\item I worked on simplifying the algebra required to derive a coherence criterion for a two patch SIS system (where dispersal matrix does not have equal coupling).
\end{itemize}

\section*{Thursday 22 March 2018}

\textbf{Solo work} \\
\emph{Approximate Duration: 1.5 Hours}
\begin{itemize}
\item Attempted to find coherence criterion for continuous time version of our SIR model (first examined the continuous time version of the SIS model).
\item Once again assumed that dispersal matrix did not have equal coupling.
\item Algebra became impossible to simplify; looking like analysis may not be as easy we thought. Only results obtained so far come from basic two patch SIS models.
\end{itemize}

\section*{Friday 23 March 2018}

\textbf{Group Meeting (in class)} \\
\emph{Approximate Duration: 1 Hour}

\begin{itemize}
\item Split up work for draft.
\item Decided that DPark and I will be doing numerical work now, while Aurora and Michael will be doing analytical work.
\end{itemize}

\section*{Sunday 25 March 2018}

\textbf{Solo work} \\
\emph{Approximate Duration: 1.5 Hours}
\begin{itemize}
\item Learned basics of Rcpp so that I can write a source file in Rcpp for a stochastic SIR model.
\item Began to tinker with Rcpp source file for our simple SIR model (no alterations); will write code tomorrow
\end{itemize}

\section*{Monday 26 March 2018}

\textbf{Group Meeting (in class)} \\
\emph{Approximate Duration: 1 Hour}
\begin{itemize}
\item Discussed direction for the paper and further split up work.
\item Analytical work will be shelved, for now; only focusing on numerical simulations.
\item I will be working on the investigation of effects of seasonal forcing and the changes in seasonal amplitude.
\end{itemize}

\textbf{Solo Work} \\
\emph{Approximate Duration: 2.5 Hours}
\begin{itemize}
\item Wrote a function in Rcpp that simulates changes for a stochastic version of our SIR model.
\item This file was then used by DPark and he copied the important bits into the "SIRmodelnpatch.cpp" file, in order to keep repository clean and code easy to work with.
\item Began to analyze different functions for $\beta (t)$. Will use sinusoidally-forced transmission rate, will also use a term-time transmission rate similar to one seen in (He et al. 2009).
\end{itemize}


\textbf{Group Meeting (evening)} \\
\emph{Approximate Duration: 1 Hour}
\begin{itemize}
\item Discussed plan for the draft.
\end{itemize}

\section*{Tuesday 27 March 2018}

\textbf{Solo work} \\
\emph{Apprxomiate Duration: 3.5 Hours}
\begin{itemize}
\item Wrote a function for term-time forced transmission rate and added it to the general source file ("SIRmodelnpatch.cpp").
\item Encountered several errors; prevalence vs. time graph looked extremely incorrect.
\item Also simulated various runs of sinusoidally-forced transmission rate when seasonal amplitude is changed, observed trends.
\item Wrote a paragraph in introduction.
\end{itemize}

\textbf{Group meeting} \\
\emph{Approximate Duration: 1 Hour}
\begin{itemize}
\item DPark and Michael helped fix the bug in my code, turned out it was a simple int/float conflict error in Rcpp.
\item Agreed that everyone will have things done by 9 AM tomorrow.
\end{itemize}

\textbf{Solo work} \\
\emph{Apprxomiate Duration: 3 Hours}
\begin{itemize}
\item Adapted the code previously used by Aurora and Michael to account for changes in seasonal amplitude (and ignore changes in $\mathcal{R}_0$ and connectivity matrices).
\item Ran simulations for different seasonal amplitude values for both forced transmission rates. Coherence trends are difficult to detect right now. A bifurcation diagram will be necessary to better understand the effects of seasonal amplitude.
\item Wrote up my findings in the project document, while also explaining a bit about time-dependent transmission rates in general.
\end{itemize}

\section*{Total time spent on this project}

\begin{quote}

\emph{Group work:} 8.5 hours

\emph{Solo work:} 16 hours

\end{quote}

\end{document}
